\documentclass[11pt]{article}
\usepackage[a4paper,margin=1in]{geometry}
\usepackage{amsmath,amssymb,amsthm}
\usepackage[T1]{fontenc}
\usepackage{lmodern}
\usepackage[hidelinks]{hyperref}
\setlength{\emergencystretch}{3em}

\theoremstyle{plain}
\newtheorem{theorem}{Theorem}
\newtheorem{lemma}[theorem]{Lemma}
\newtheorem{proposition}[theorem]{Proposition}
\newtheorem{corollary}[theorem]{Corollary}
\theoremstyle{remark}
\newtheorem{remark}[theorem]{Remark}

\DeclareMathOperator{\Gal}{Gal}

\title{A Counterexample to Conjecture 00000000269:\\
the Galois Group of $x^n+x+1$ Is Not $S_n$ When $n\equiv2\pmod 3$}
\author{%
  Feiyu\thanks{Independent researcher.}
  \and
  Claude (Opus 5.5)\thanks{Anthropic.}%
}
\date{2026-10-02}

\begin{document}
\maketitle

\begin{abstract}
Conjecture 00000000269 asserts that for all sufficiently large $n$ the Galois
group of $x^n+x+1$ is the symmetric group $S_n$. It is false: for every
$n\equiv2\pmod3$ the trinomial is divisible by $x^2+x+1$, and its Galois group
then has at most $2\,(n-2)!<n!$ elements, so it is not even abstractly
isomorphic to $S_n$. Such $n$ are unbounded, so no threshold works. The
refutation is formalized in Lean~4 against Mathlib.
\end{abstract}

\section{The conjecture as stated}

The original text of conjecture 00000000269 reads:

\begin{quote}
Definition: The Galois group of $x^n + x + 1$. Conjecture: For all
sufficiently large $n$, this group is $S_n$ (Selmer proved the corresponding
statement for $x^n - x - 1$).
\end{quote}

\section{Reading of the statement}

The polynomial $f_n(x)=x^n+x+1$ has integer coefficients, and its Galois group
is taken over $\mathbb Q$, as in the work of Selmer to which the statement
refers. $\Gal(f_n)$ denotes the automorphism group of a splitting field of
$f_n$ over $\mathbb Q$.

``This group is $S_n$'' is usually understood as: $\Gal(f_n)$, acting on the $n$
roots of $f_n$, is the full symmetric group on them. Any such reading implies
the weakest one,
\begin{equation}\label{eq:weak}
  \Gal(f_n)\cong S_n\quad\text{as abstract groups},
\end{equation}
and this is the statement we refute. So the conjecture asserts that there is an
$N$ such that \eqref{eq:weak} holds for every $n\ge N$, and we show that
\eqref{eq:weak} fails for infinitely many $n$.

\section{Result}

\begin{theorem}\label{thm:main}
Let $n\ge5$ with $n\equiv2\pmod 3$. Then $|\Gal(f_n)|\le 2\,(n-2)!<n!$. In
particular $\Gal(f_n)\not\cong S_n$, and Conjecture 00000000269 is false.
\end{theorem}

\section{Proof}

\begin{lemma}\label{lem:dvd}
If $n=3k+5$ with $k\ge0$, then $x^2+x+1$ divides $x^n+x+1$ in $\mathbb Q[x]$.
\end{lemma}

\begin{proof}
We have the identity
\[
  x^{3k+5}+x+1 \;=\; x^2\bigl((x^3)^{k+1}-1\bigr)+\bigl(x^2+x+1\bigr).
\]
Since $x^3-1$ divides $(x^3)^{k+1}-1$ and $x^3-1=(x-1)(x^2+x+1)$, both terms on
the right are divisible by $x^2+x+1$. (Equivalently: a primitive cube root of
unity $\omega$ satisfies $\omega^n+\omega+1=\omega^2+\omega+1=0$.)
\end{proof}

\begin{lemma}\label{lem:deg}
If $p\in\mathbb Q[x]$ has degree $d$, then $|\Gal(p)|\le d!$.
\end{lemma}

\begin{proof}
$\Gal(p)$ permutes the roots of $p$ in its splitting field $L$, and this action
is faithful because $L$ is generated over $\mathbb Q$ by those roots. There are
at most $d$ roots, so $\Gal(p)$ embeds in a symmetric group on at most $d$
letters.
\end{proof}

\begin{lemma}\label{lem:prod}
For $p,q\in\mathbb Q[x]$, $|\Gal(pq)|\le|\Gal(p)|\cdot|\Gal(q)|$.
\end{lemma}

\begin{proof}
The splitting field $M$ of $pq$ contains splitting fields of $p$ and of $q$,
and every automorphism of $M$ maps each of them to itself. Restriction
therefore gives a homomorphism $\Gal(pq)\to\Gal(p)\times\Gal(q)$. It is
injective, because $M$ is generated by the roots of $p$ and of $q$, so an
automorphism of $M$ that fixes both subfields is the identity.
\end{proof}

\begin{proof}[Proof of Theorem~\ref{thm:main}]
Write $n=3k+5$. By Lemma~\ref{lem:dvd}, $f_n=(x^2+x+1)\,g$ for some
$g\in\mathbb Q[x]$, and comparing degrees gives $\deg g=n-2$. By
Lemmas~\ref{lem:prod} and~\ref{lem:deg},
\[
  |\Gal(f_n)|\;\le\;|\Gal(x^2+x+1)|\cdot|\Gal(g)|\;\le\;2!\,(n-2)!\;=\;2\,(n-2)!.
\]
On the other hand $n!=n(n-1)\,(n-2)!$ and $n(n-1)\ge20>2$, so
$|\Gal(f_n)|<n!=|S_n|$ and \eqref{eq:weak} fails.

Given any threshold $N$, the integer $n=3N+5$ satisfies $n\ge N$, $n\ge5$ and
$n\equiv2\pmod3$, so \eqref{eq:weak} fails for it. Hence no threshold has the
property asserted by the conjecture.
\end{proof}

\section{Formalization}

The accompanying Lean~4 project (\texttt{lean/Submission/Basic.lean}) states
the conjecture and proves its negation against Mathlib, whose
\texttt{Polynomial.Gal p} is the Galois group of the splitting field of $p$.

\begin{itemize}
  \item \texttt{trinomial n} --- the polynomial $x^n+x+1\in\mathbb Q[x]$.
  \item \texttt{GalIsSymmetric n} --- \eqref{eq:weak}: a group isomorphism
    \texttt{(trinomial n).Gal}~$\simeq^{*}$~\texttt{Equiv.Perm (Fin n)}
    exists.
  \item \texttt{ConjectureHolds} --- there is $N$ with
    \texttt{GalIsSymmetric n} for all $n\ge N$.
  \item \texttt{card\_gal\_le} --- Lemma~\ref{lem:deg}, from Mathlib's
    \texttt{Gal.galActionHom\_injective} and
    \texttt{ncard\_rootSet\_le}.
  \item \texttt{card\_gal\_mul\_le} --- Lemma~\ref{lem:prod}, from Mathlib's
    \texttt{Gal.restrictProd\_injective}.
  \item \texttt{cyclotomic\_dvd\_trinomial} --- Lemma~\ref{lem:dvd};
    \texttt{natDegree\_cyclotomic}, \texttt{natDegree\_trinomial} --- the
    degrees $2$ and $n$.
  \item \texttt{card\_gal\_trinomial\_le} --- $|\Gal(f_{3k+5})|\le 2\,(3k+3)!$;
    \texttt{not\_galIsSymmetric} --- \eqref{eq:weak} fails for $n=3k+5$.
  \item \texttt{conjecture\_00000000269\_false} ---
    $\lnot\,$\texttt{ConjectureHolds}, by taking $n=3N+5$.
\end{itemize}

The project builds with \texttt{lake build} on
\texttt{leanprover/lean4:v4.33.1} with Mathlib \texttt{v4.33.1}. It contains
no \texttt{sorry}, no \texttt{admit} and no \texttt{native\_decide}, and
introduces no axioms: \texttt{\#print axioms
conjecture\_00000000269\_false} reports exactly \texttt{propext},
\texttt{Classical.choice} and \texttt{Quot.sound}.

\section{Remarks}

The divisibility in Lemma~\ref{lem:dvd} is classical. Selmer \cite{selmer}
showed that $x^n+x+1$ is irreducible over $\mathbb Q$ for $n\not\equiv2\pmod
3$, and that for $n\equiv2\pmod3$, $n>2$, it factors as $(x^2+x+1)\,g_n(x)$
with $g_n$ irreducible. So the obstruction found here is the only reducibility
of the family, and it occurs for one residue class in three. The conjecture
could be repaired by restricting to $n\not\equiv2\pmod3$, or by replacing
$x^n+x+1$ with $g_n$ when $n\equiv2\pmod3$; we do not address either version.
The statement Selmer proved for $x^n-x-1$, which the conjecture cites, is
irreducibility for every $n$. That polynomial has no such cyclotomic factor.

\begin{thebibliography}{9}
\bibitem{selmer} E.~S.~Selmer, \emph{On the irreducibility of certain
  trinomials}, Math. Scand. 4 (1956), 287--302.
\end{thebibliography}

\end{document}
